\begin{table}[H]
  \centering
  \caption{The statements of Part~I and what each rests on.}
  \label{tab:foreclosure:summary}
  \small
  \renewcommand{\arraystretch}{1.1}
  \begin{tabularx}{\textwidth}{@{}>{\raggedright\arraybackslash}p{0.25\textwidth}>{\raggedright\arraybackslash}X>{\raggedright\arraybackslash}p{0.27\textwidth}@{}}
    \toprule
    Statement & Content & Status \\
    \midrule
    \Cref{thm:carrier-dichotomy} & native-outcome branch for \nonXequiv\ carriers; exactly one failure state for \Xequiv\ ones & proved, assuming (Cls) \\
    \Cref{cor:bridge-impossibility} & a bridge into an \Xequiv\ composite carries an \XstrengthObl & proved, assuming (Cls) and (BO) \\
    \Cref{thm:coverage} & the four states are exhaustive & schema; follows from (Cls) when some claim expresses \(R_{C}=0\) \\
    \Cref{thm:type-finite-foreclosure} & audit outcomes form a four-element set & proved, assuming coverage \\
    \Cref{thm:attack-foreclosure-v4} & internal moves do not advance an in-scope target without outside content & schema \\
    \Cref{thm:ctmt-recursion} & matrix-element termination nests to any finite depth & proved \\
    \bottomrule
  \end{tabularx}
\end{table}
